%% kapa: exam=2 round=25 year=2014 period=2 subject=appraisal_theory subject_ko=감정평가이론 seq=1 number=1 points=40
%% source: Q-Net 감정평가사 기출문제 (cst00309) artlSeq=5109209 file=extracted/1008/5109209/2014년 제25회 감정평가사 2차시험 기출문제/2교시(감정평가이론).hwp
%% conversion: kordoc 4.24.51 -> tools/gen2.py (auto); review=proofread-sonnet; flags=none
\begin{longquestion}{【문제 1】}
최근 부동산시장 환경변화로 부동산감정평가에서 고려할 사항이 늘고 있다. 감정평가 원리 및 방식에 대한 다음 물음에 답하시오. (40점) \par
1) 리모델링된 부동산에 대해 감정평가3방식을 적용하여 감정평가할 때 유의할 사항을 설명하시오. (10점) \par
2) 토양오염이 의심되는 토지에 대한 감정평가안건의 처리방법을 설명하시오. (15점) \par
3) 공익사업을 위해 수용될 지구에 포함되어 장기 미사용 중이던 토지가 해당 공익사업의 중단으로 지구지정이 해제되었을 때, 당해 토지 및 주변부 토지에서 초래될 수 있는 경제적 손실을 부동산평가원리에 근거하여 설명하시오. (15점) \par
\end{longquestion}
